\hypertarget{appendix-390ea641da54}{}
\section*{Machine epsilon}
\section{Machine epsilon}\label{machine-epsilon}

\textbf{Curriculum depth:} Introduced · \textbf{First seen:} C004 · \textbf{Latest development:} C004

\subsection{Definition and mechanism}\label{definition-and-mechanism}

For these binary types epsilon is upward spacing at 1: 2\^{}-23 for binary32 and 2\^{}-52 for binary64. It is neither minimum positive value nor a physical threshold.

\subsection{Prerequisites}\label{prerequisites}

\begin{itemize}
\tightlist
\item
  \hyperref[appendix-a24fd710cf67]{Significand and hidden bit} --- Fraction weights and effective precision determine value and the smallest retained position.
\end{itemize}

\subsection{Used by}\label{used-by}

\begin{itemize}
\tightlist
\item
  \hyperref[appendix-14bab4cab9d4]{ULP and representable spacing} --- uses this mechanism in its explanation.
\item
  \hyperref[appendix-0ec599f8ba27]{Numerical tolerance policies} --- uses this mechanism in its explanation.
\end{itemize}

\subsection{Recall and next depth}\label{recall-and-next-depth}

Reconstruct the concrete example in the definition and identify its validity assumptions.

\textbf{Pending:} Explain min, denorm\_min, lowest, and task-specific tolerance independently.

\subsection{Source and navigation}\label{source-and-navigation}

\href{Continuous\%20Notes/Computing\%20Foundations\%20-\%20From\%20Bits\%20to\%20C++\%20and\%20Systems\%20APIs\#A\%20stretching\%20ruler\%20for\%20real-valued\%20quantities}{M004 teaching chapter} · \hyperref[appendix-fd090d982a15]{Per-part definitions} · \href{}{Dependency Map (02-Foundations-II, appendix-4ad47bdd269a)}

\subsection{Epsilon is local to one magnitude}\label{epsilon-is-local-to-one-magnitude}

For a floating type, machine epsilon is commonly defined as the distance from \(1\) to the next representable number greater than \(1\). For binary precision \(p\), that spacing is \(2^{1-p}\). It describes one location in the representable lattice; it is \textbf{not} a universal comparison tolerance.

Floating spacing scales with exponent. Near a large world coordinate, one ULP can be much larger than epsilon; near zero, subnormal spacing follows a different regime. Therefore testing

\texttt{abs(a-b)\ \textless{}\ epsilon}

for arbitrary geometric or physical quantities is usually unjustified.

A defensible comparison begins with the domain: absolute tolerance for quantities near zero, relative tolerance when scale matters, ULP bounds when representation-level proximity is intended, or an exact predicate when the algorithm can provide one. Tests should include values near 1, large magnitudes, tiny normals/subnormals and adjacent representable values using \texttt{nextafter}.
